% !TeX program = pdflatex
%==============================================================================
%  Calibrated Cross-Encoder Micro-Decisions for RAG Pipelines  -- REVISION 1
%  ---------------------------------------------------------------------------
%  AUTHOR NOTE (delete before submission):
%   * The submitted version is kept in paper_revision/submitted_v1.tex.
%   * Every table body and every number in the prose is GENERATED from the
%     result files by  python scripts/make_tables.py  into paper_revision/tables/.
%     Do not type numbers by hand: re-run the script after each experiment.
%   * Anything not yet measured shows as a RED \ph{--}. Before submission,
%       grep -c "ph{" paper_revision/tables/*.tex   must print 0 everywhere.
%   * BLUE [\revnote{...}] marks prose whose wording depends on a result that
%     is not in yet. Re-read and adjust it, then delete the note.
%   * Camera-ready: \renewcommand{\ph}[1]{#1} and \renewcommand{\revnote}[1]{}.
%==============================================================================

\documentclass[twocolumn]{ICCKjournal}

\usepackage{amsmath}
\usepackage{amssymb}
\usepackage{amsthm}
\usepackage[none]{hyphenat}
\usepackage{enumitem}
\usepackage{float}
\usepackage{booktabs}
\usepackage{multirow}
\usepackage{subcaption}
\usepackage{xcolor}
\usepackage{algorithm}
\usepackage{algpseudocode}
\usepackage{adjustbox}
\usepackage{xurl}

\graphicspath{{images/},{figures/},{res/},{media/}}

\AtBeginEnvironment{thebibliography}{\small}

% ---- placeholder marker: red = not yet measured ------------------------------
\newcommand{\ph}[1]{\textcolor{red}{#1}}
% ---- revision note: blue = wording to re-check once the result is in --------
\newcommand{\revnote}[1]{\textcolor{blue}{[#1]}}
% ---- all numbers quoted in the text (generated) -----------------------------
% generated by scripts/make_tables.py - do not edit by hand
\newcommand{\BGEGroHeld}{\ph{--}}
\newcommand{\BGEGroTest}{\ph{--}}
\newcommand{\BGERelHeld}{\ph{--}}
\newcommand{\BGERelTest}{\ph{--}}
\newcommand{\BGESufHeld}{\ph{--}}
\newcommand{\BGESufTest}{\ph{--}}
\newcommand{\CascGrdHeldAcc}{\ph{--}}
\newcommand{\CascGrdHeldEsc}{\ph{--}}
\newcommand{\CascGrdHeldParityT}{\ph{--}}
\newcommand{\CascGrdLLMHeldAcc}{\ph{--}}
\newcommand{\CascGrdLLMMs}{\ph{--}}
\newcommand{\CascGrdMs}{\ph{--}}
\newcommand{\CascGrdNanoHeldAcc}{\ph{--}}
\newcommand{\CascGrdSpeedup}{\ph{--}}
\newcommand{\CascGrdTestAcc}{\ph{--}}
\newcommand{\CascGrdTestEsc}{\ph{--}}
\newcommand{\CascT}{\ph{--}}
\newcommand{\CustAgnewsRand}{0.250}
\newcommand{\CustAgnewsRev}{\ph{--}}
\newcommand{\CustAgnewsVOne}{0.186}
\newcommand{\CustAgnewsVZero}{0.396}
\newcommand{\CustEmotionRand}{0.167}
\newcommand{\CustEmotionRev}{\ph{--}}
\newcommand{\CustEmotionVOne}{0.138}
\newcommand{\CustEmotionVZero}{0.334}
\newcommand{\CustYahooRand}{0.100}
\newcommand{\CustYahooRev}{\ph{--}}
\newcommand{\CustYahooVOne}{0.056}
\newcommand{\CustYahooVZero}{0.239}
\newcommand{\EEHotBgeFone}{\ph{--}}
\newcommand{\EEHotBgeMs}{\ph{--}}
\newcommand{\EEHotNJFone}{\ph{--}}
\newcommand{\EEHotNJMs}{\ph{--}}
\newcommand{\EEHotNoneFone}{\ph{--}}
\newcommand{\EEHotNoneMs}{\ph{--}}
\newcommand{\EEHotQwenFone}{\ph{--}}
\newcommand{\EEHotQwenMs}{\ph{--}}
\newcommand{\EEMusBgeFone}{\ph{--}}
\newcommand{\EEMusBgeMs}{\ph{--}}
\newcommand{\EEMusNJFone}{\ph{--}}
\newcommand{\EEMusNJMs}{\ph{--}}
\newcommand{\EEMusNoneFone}{\ph{--}}
\newcommand{\EEMusNoneMs}{\ph{--}}
\newcommand{\EEMusQwenFone}{\ph{--}}
\newcommand{\EEMusQwenMs}{\ph{--}}
\newcommand{\LLMGroHeld}{\ph{--}}
\newcommand{\LLMGroTest}{\ph{--}}
\newcommand{\LLMRelHeld}{\ph{--}}
\newcommand{\LLMRelTest}{\ph{--}}
\newcommand{\LLMSufHeld}{\ph{--}}
\newcommand{\LLMSufTest}{\ph{--}}
\newcommand{\NCalib}{3,982}
\newcommand{\NDev}{2,488}
\newcommand{\NJEceHeldMax}{\ph{--}}
\newcommand{\NJEceHeldMin}{\ph{--}}
\newcommand{\NJEceTestMax}{\ph{--}}
\newcommand{\NJEceTestMin}{\ph{--}}
\newcommand{\NJGroHeld}{\ph{--}}
\newcommand{\NJGroTest}{\ph{--}}
\newcommand{\NJMsMax}{\ph{--}}
\newcommand{\NJMsMin}{\ph{--}}
\newcommand{\NJRelHeld}{\ph{--}}
\newcommand{\NJRelTest}{\ph{--}}
\newcommand{\NJSufHeld}{\ph{--}}
\newcommand{\NJSufTest}{\ph{--}}
\newcommand{\NJTauGro}{\ph{--}}
\newcommand{\NJTauRel}{\ph{--}}
\newcommand{\NJTauSuf}{\ph{--}}
\newcommand{\NSeeds}{\ph{--}}
\newcommand{\NTest}{3,976}
\newcommand{\NTestGrd}{500}
\newcommand{\NTestRel}{1,984}
\newcommand{\NTestSuf}{1,492}
\newcommand{\NTrain}{47,272}
\newcommand{\OPGroHeld}{\ph{--}}
\newcommand{\OPGroTest}{\ph{--}}
\newcommand{\OPRelHeld}{\ph{--}}
\newcommand{\OPRelTest}{\ph{--}}
\newcommand{\OPSufHeld}{\ph{--}}
\newcommand{\OPSufTest}{\ph{--}}
\newcommand{\SpeedupMax}{\ph{--}}
\newcommand{\SpeedupMin}{\ph{--}}
\newcommand{\TruncInDistMax}{\ph{--}}
\newcommand{\TruncMusSufAccBal}{\ph{--}}
\newcommand{\TruncMusSufAccNT}{\ph{--}}
\newcommand{\TruncMusSufAccT}{\ph{--}}
\newcommand{\TruncMusSufGold}{\ph{--}}
\newcommand{\TruncMusSufPct}{\ph{--}}

% ---- notation ---------------------------------------------------------------
\newcommand{\logit}{\ell}
\newtheorem{proposition}{Proposition}
\newtheorem{remark}{Remark}

\articletype{Research Article}
\editor{Editor}
\submitdate{2026}
\acceptdate{2026}
\pubdate{2026}

\jnlsetup{%
    author  = {Divakar, S. \& Priyadarshini, R.},
    journal = {ICCK Transactions and Journals},
    title   = {CaliDec: Calibrated Cross-Encoder Micro-Decisions for Retrieval-Augmented Generation with Confidence Cascading},
    ctitle  = {Calibrated Cross-Encoder Micro-Decisions for RAG},
    volume  = {},
    number  = {},
    year    = {2026},
    doi     = {}
}

\dataavail{%
The complete implementation (model, training pipeline, evaluation and baseline scripts, calibration code, statistical tests and unit tests) is publicly available at \url{https://github.com/shubham10divakar/nano-jev}. Pre-trained weights are hosted at \url{https://huggingface.co/sdmlai/nano-jev}. All training and evaluation datasets are publicly available from their original sources.}

\setcounter{page}{1}
\makeatother
\begin{document}

\author[1][shubham.divakar@gmail.com]{Subham Divakar}[0000-0003-0635-1058]
\author[2][priyadarshini.rojalina@gmail.com]{Rojalina Priyadarshini}[0000-0002-5481-5251]

\affil[1]{Department of Electronics and Telecommunications, C.V. Raman Global University, Bhubaneswar, India}
\affil[2]{Department of Computer Science, C.V. Raman Global University, Bhubaneswar, India}

\ifproof
\input{proof}
\fi

\maketitle
\abstract{Retrieval-Augmented Generation (RAG) pipelines make many small, repeated decisions at inference time---is this passage relevant? is the retrieved context sufficient? is the generated answer grounded?---yet most systems delegate these judgments either to the same large language model (LLM) that generates the answer, at high latency and cost, or to heuristic score thresholds that lack probabilistic meaning. We introduce \textbf{CaliDec}, a 33.4\,M-parameter cross-encoder \emph{micro-decision} model that casts each RAG judgment as a typed, closed-set classification problem and returns \emph{temperature-calibrated} probabilities over the decision's options. Each (question, option, context) triple is scored independently by a single-logit cross-encoder head, which makes the output invariant to option order; per-decision temperatures are fitted on a calibration split that is used for nothing else. Three built-in decision types are supported: \emph{relevance} (3-way), \emph{sufficiency} (binary) and \emph{groundedness} (binary). Averaged over \NSeeds{} training seeds, CaliDec reaches \NJRelTest{}, \NJSufTest{} and \NJGroTest{} test accuracy on held-out halves of the HotpotQA, SQuAD~2.0 and MultiNLI validation sets, with expected calibration error \NJEceTestMin--\NJEceTestMax{}, at \NJMsMin--\NJMsMax{}\,ms per decision on a consumer GPU (\SpeedupMin--\SpeedupMax$\times$ faster than a prompted 4\,B-parameter LLM). We compare against a fine-tuned single-pass classifier on the same backbone and data, compact groundedness checkers (NLI, AlignScore, MiniCheck), a 568\,M-parameter reranker and the prompted LLM, with paired significance tests. On unseen datasets (MuSiQue, VitaminC, LLM-AggreFact) every model loses accuracy, and temperatures fitted in-distribution become overconfident (ECE \NJEceHeldMin--\NJEceHeldMax). We further evaluate a \emph{confidence cascade} that escalates low-confidence decisions to the LLM on all three decision types: with the threshold chosen on the calibration split, the cascade reaches \CascGrdTestAcc{} on in-distribution groundedness while escalating \CascGrdTestEsc{} of decisions, and it automatically escalates (almost) nothing where the LLM is the weaker model. A small end-to-end RAG evaluation measures the downstream effect of using CaliDec as a passage filter. The model, training pipeline, evaluation code and weights are released as open-source software. \revnote{Re-read every claim in this abstract against the final tables; if the 1-pass classifier or MiniCheck matches CaliDec, lead with the typed interface, calibration and cascade rather than accuracy.}}

\keywords{Retrieval-augmented generation, cross-encoder, decision model, calibration, confidence cascade, natural language inference, question answering.}


\begin{figure}[b]
	\citationblock
	\vspace*{-.98em}
\end{figure}

%═══════════════════════════════════════════════════════════════════════════════
\section{Introduction}
\label{sec:intro}
%═══════════════════════════════════════════════════════════════════════════════

Retrieval-Augmented Generation (RAG)~\cite{lewis2020rag} has become a standard approach for grounding large language model (LLM) outputs in external knowledge. A typical RAG pipeline retrieves passages from a corpus, determines which are useful, decides whether the retrieved context suffices, generates an answer, and then checks whether that answer is faithful to the evidence. Each of these intermediate steps requires a \emph{judgment}---a small decision with a fixed set of outcomes.

In practice, these decisions are handled in one of three ways, none of which is fully satisfactory. (i)~The generating LLM is prompted to self-assess, which is expensive, latency-bound, and produces poorly calibrated confidences~\cite{kadavath2022language, kuhn2023semantic}. (ii)~A reranker or NLI model provides a single relevance or entailment score, which must be mapped post hoc to the decision's categories via hand-tuned thresholds~\cite{bge_m3}. (iii)~The pipeline uses deterministic heuristics (e.g.\ BM25 score $>$ 0.5), which carry no uncertainty information and cannot be combined meaningfully.

We argue that these judgments deserve a dedicated model class: small, fast, \emph{typed} decision models that accept a question, a closed set of options, and a context, and return calibrated probabilities. The key properties are: (i)~each decision type has an explicit schema of options, so the output space is always well-defined; (ii)~the model never generates text, eliminating parsing failures and hallucinated self-assessments; (iii)~the probabilities are temperature-calibrated, making thresholds frequency-interpretable and enabling principled cascading to a larger model when the small model is uncertain.

\textbf{Contributions.}
\begin{enumerate}[leftmargin=*]
    \item \textbf{Typed decision formulation.} We formalise the recurring micro-decisions of RAG pipelines as closed-set classification problems with a shared cross-encoder that scores each option separately. Three built-in types---relevance, sufficiency and groundedness---cover the most common judgments (Section~\ref{sec:decisions}).
    \item \textbf{CaliDec model and a controlled ablation.} We present a 33.4\,M-parameter K-pass cross-encoder with per-decision temperature scaling (Section~\ref{sec:model}) and compare it with a standard fine-tuned 1-pass classifier on the same backbone, data and protocol (Section~\ref{sec:onepass}), so that the effect of the K-pass design is isolated. \revnote{State the outcome in one clause once Table~\ref{tab:onepass} is in.}
    \item \textbf{Calibration under a clean protocol.} Checkpoints are selected on a development split carved from the training pool; temperatures, cascade thresholds and baseline mappings are fitted on a calibration split used for nothing else. We report seed variance, bootstrap confidence intervals and paired McNemar tests, and quantify how in-distribution calibration degrades on unseen domains (Section~\ref{sec:calresults}).
    \item \textbf{Confidence cascade on all decisions.} We evaluate a cascade that escalates uncertain decisions to a 4\,B-parameter LLM for all three decision types, on in-distribution and held-out data, and show that a threshold chosen on the calibration split avoids escalating to a model that is weaker for that decision (Section~\ref{sec:cascaderesults}).
    \item \textbf{RAG-oriented evaluation.} Beyond the decision benchmarks, we measure input truncation, groundedness of multi-sentence LLM-generated claims (LLM-AggreFact) and a small end-to-end RAG pipeline (Sections~\ref{sec:truncation}--\ref{sec:e2e}). Code and weights are released.
\end{enumerate}

The remainder of the paper is organised as follows. Section~\ref{sec:related} reviews related work. Section~\ref{sec:method} describes the model, decision schema and calibration. Section~\ref{sec:experiments} details the experimental protocol, Section~\ref{sec:results} presents and discusses results, and Section~\ref{sec:conclusion} concludes.

%═══════════════════════════════════════════════════════════════════════════════
\section{Related Work}
\label{sec:related}
%═══════════════════════════════════════════════════════════════════════════════

\subsection{Retrieval-Augmented Generation}
RAG~\cite{lewis2020rag} augments LLM generation with retrieved passages. REALM~\cite{guu2020realm} jointly trains retriever and reader. Self-RAG~\cite{asai2024selfrag} adds reflection tokens so the LLM critiques its own retrieval and generation. Corrective RAG (CRAG)~\cite{yan2024corrective} introduces a lightweight evaluator that triages retrieved documents before generation. Our work shares the insight that retrieval evaluation should be separated from generation, but focuses on a dedicated, typed, calibrated decision model rather than special tokens or rule-based evaluators.

\subsection{Cross-Encoders for Relevance and NLI}
Cross-encoders~\cite{nogueira2019passage} process query--document pairs jointly and outperform bi-encoders on reranking tasks. Models such as the MS MARCO MiniLM cross-encoder~\cite{reimers2019sentence} and BGE reranker~\cite{bge_m3} are widely used for passage reranking. For NLI, BERT~\cite{devlin2019bert}, RoBERTa~\cite{liu2019roberta} and DeBERTa-v3~\cite{he2023debertav3} fine-tuned on MultiNLI, FEVER and ANLI~\cite{laurer2024less} remain strong baselines. Our approach extends cross-encoders from a single relevance score to multi-type decisions with calibrated option probabilities.

\subsection{Factual Consistency and RAG Evaluation}
A large body of work checks whether a text is supported by a source document. SummaC~\cite{laban2022summac} applies sentence-level NLI and aggregates over document chunks; TRUE~\cite{honovich2022true} standardises factual-consistency benchmarks and shows that NLI-based checkers are strong; AlignScore~\cite{zha2023alignscore} trains a unified alignment function on many tasks; MiniCheck~\cite{tang2024minicheck} trains small checkers on synthetic data and introduces LLM-AggreFact, on which it matches far larger LLMs. For RAG specifically, RAGAS~\cite{es2024ragas} scores faithfulness and relevance with LLM prompts, and ARES~\cite{saadfalcon2024ares} fine-tunes lightweight judges for context relevance, answer faithfulness and answer relevance and corrects their estimates with a small labelled set. CaliDec differs from these in three ways: (i)~one model serves several \emph{typed} decisions with explicit option sets (including the 3-way relevance and the pre-generation sufficiency decision, which consistency checkers do not cover); (ii)~its outputs are temperature-calibrated per decision so that thresholds are meaningful at run time, not only for aggregate evaluation; and (iii)~it is designed as the first stage of a confidence cascade. Groundedness checkers are therefore direct competitors on one of our three decisions, and we compare against AlignScore and MiniCheck there (Section~\ref{sec:groundedbaselines}).

\subsection{Confidence Estimation and Calibration}
Modern neural networks are often poorly calibrated~\cite{guo2017calibration}. Temperature scaling~\cite{guo2017calibration} is the simplest post-hoc calibration method, fitting a single scalar per task. Platt scaling~\cite{platt1999probabilistic} generalises this with a logistic regression. Kadavath et al.~\cite{kadavath2022language} study the self-evaluation capabilities of LLMs, finding that larger models are better calibrated but still unreliable on adversarial inputs. We apply per-decision temperature scaling and show it is effective in-distribution but insufficient for domain transfer.

\subsection{Cascading and Adaptive Inference}
FrugalGPT~\cite{chen2023frugalgpt} cascades through progressively larger LLMs based on confidence, reducing API cost by up to 98\%. Routing-based approaches~\cite{ding2024hybrid} select between local and cloud models per query. Speculative decoding~\cite{leviathan2023fast} uses a small model to draft and a large model to verify. Our cascade is simpler: a single-model confidence threshold decides whether to accept or escalate, enabled by the calibrated probabilities that make the threshold meaningful.

\subsection{The Jev Decision-Model Paradigm}
The Jev architecture~\cite{jpalmer2026jev} (TypeSafe AI) introduced the concept of a compact, typed, non-generative ``System~1'' decision model for agentic workflows. Jev demonstrated that many runtime judgments---routing, relevance, sufficiency---can be handled by a small classifier instead of an LLM, with dramatically lower latency. Our work is an independent, open-source realisation of this paradigm, with the additions of explicit probability calibration, a formal cascade strategy, and systematic evaluation on held-out domains.

%═══════════════════════════════════════════════════════════════════════════════
\section{Proposed Method}
\label{sec:method}
%═══════════════════════════════════════════════════════════════════════════════

\subsection{Overview}
Fig.~\ref{fig:architecture} illustrates the CaliDec inference pipeline. A typed decision specifies a question template, a set of options, and a context. For each option, a cross-encoder produces one logit. The per-option logits are divided by a decision-specific temperature $\tau_d$ and softmaxed into calibrated probabilities. In the cascade variant, if the top probability falls below a confidence threshold $t$, the decision is forwarded to a large language model.

\begin{figure*}[!t]
    \centering
    \includegraphics[width=0.92\textwidth]{CaliDec_Architecture.png}
    \caption{CaliDec inference pipeline. A typed decision (e.g.\ \emph{relevance}) specifies a question template and option set. Each (question + option, context) pair is encoded by a cross-encoder that outputs one logit. The option logits are temperature-scaled and softmaxed. If the cascade is active, decisions with top probability below threshold $t$ are escalated to an LLM.}
    \label{fig:architecture}
\end{figure*}

\subsection{Typed Decision Schema}
\label{sec:decisions}
We define a \emph{decision type} $d$ as a triple $(\mathcal{Q}_d, \mathcal{O}_d, \mathcal{S}_d)$: a question template $\mathcal{Q}_d$ parameterised by task-specific fields, a fixed set of options $\mathcal{O}_d = \{o_1, \dots, o_K\}$, and a context type $\mathcal{S}_d$. Table~\ref{tab:decisions} lists the three built-in types.

\begin{table}[H]
\caption{Built-in decision types. Each type has a question template, a fixed option set, and a context specification.}
\label{tab:decisions}
\centering
\small
\begin{adjustbox}{max width=\linewidth}
\begin{tabular}{llll}
\toprule
\textbf{Type} & \textbf{Options} & \textbf{Question template} & \textbf{Context} \\
\midrule
relevance & \{irrelevant, partially & How relevant is this passage & single \\
 & relevant, directly answers\} & to the query: \texttt{\{query\}} & passage \\
\addlinespace
sufficient & \{yes, no\} & Does the context contain enough & multiple \\
 & & information to answer: \texttt{\{query\}} & passages \\
\addlinespace
grounded & \{yes, no\} & Is this claim supported by the & evidence \\
 & & context? Claim: \texttt{\{claim\}} & text \\
\bottomrule
\end{tabular}
\end{adjustbox}
\end{table}

The interface also accepts a custom \texttt{decide} call with any question and option set, since options are scored independently (Section~\ref{sec:model}). We stress that this is an \emph{interface} property, not a demonstrated capability: Section~\ref{sec:custom} measures zero-shot accuracy on unseen label sets and finds it unreliable for the released models, so we restrict our claims to the three trained decision types.

\subsection{Cross-Encoder Scoring}
\label{sec:model}
For a decision instance with question $q$, options $\{o_1, \dots, o_K\}$ and context $s$, CaliDec constructs $K$ input pairs. The $k$-th pair is:
\begin{equation}
    x_k = \texttt{[CLS]}\ \text{question: } q \text{ option: } o_k\ \texttt{[SEP]}\ s\ \texttt{[SEP]}
    \label{eq:input}
\end{equation}
Each $x_k$ is processed by a cross-encoder $f_\theta$ (a transformer encoder followed by a single-logit linear head) to produce a scalar logit $\logit_k = f_\theta(x_k) \in \mathbb{R}$. The raw probability distribution over options is:
\begin{equation}
    p_k = \frac{\exp(\logit_k / \tau_d)}{\sum_{j=1}^{K} \exp(\logit_j / \tau_d)},
    \label{eq:softmax}
\end{equation}
where $\tau_d$ is the decision-specific temperature (Section~\ref{sec:calibration}).

\textbf{Independence property.} Because each option is scored in a separate forward pass, the logits and hence the final probabilities are invariant to option ordering, and adding or removing an option changes only the softmax normalisation. This costs $K$ forward passes per decision instead of one; Section~\ref{sec:onepass} measures what the design gains and costs against a conventional 1-pass classifier.

\textbf{Architecture.} CaliDec v1.0 uses \texttt{microsoft/MiniLM-L12-H384-uncased}~\cite{wang2020minilm} as the base encoder: 12 transformer layers, hidden size 384, with a single-output classification head. The total parameter count is 33.4\,M. The maximum input length is 512 tokens; context exceeding this limit is truncated from the right (\texttt{truncation="only\_second"}); Section~\ref{sec:truncation} measures how often this happens.

\textbf{Training objective.} Given an example with $K$ options and ground-truth label $y \in \{1, \dots, K\}$, the training loss is the cross-entropy over the grouped logits:
\begin{equation}
    \mathcal{L} = -\log \frac{\exp(\logit_y)}{\sum_{k=1}^{K} \exp(\logit_k)}.
    \label{eq:loss}
\end{equation}
Because the logits are grouped per example (not per pair), the model learns to rank the correct option above alternatives, not merely to produce high absolute scores.

\subsection{Temperature Calibration}
\label{sec:calibration}
After training, we fit a per-decision temperature $\tau_d > 0$ that minimises the negative log-likelihood on the calibration split:
\begin{equation}
    \tau_d^\star = \arg\min_{\tau > 0}\; -\frac{1}{N_{\mathrm{cal}}} \sum_{i=1}^{N_{\mathrm{cal}}} \log p_{y_i}(\tau),
    \label{eq:tempfit}
\end{equation}
where $p_{y_i}(\tau) = \mathrm{softmax}(\boldsymbol\logit_i / \tau)_{y_i}$. We optimise $\log\tau$ with L-BFGS~\cite{liu1989lbfgs} in the range $[\exp(-3), \exp(3)]$ and fall back to $\tau=1$ if calibration does not improve NLL. The calibration split is not used for training or checkpoint selection (Section~\ref{sec:protocol}). The fitted temperatures (mean over seeds) are relevance $\tau = \NJTauRel$, sufficient $\tau = \NJTauSuf$ and grounded $\tau = \NJTauGro$; values above 1 indicate that the raw model is overconfident. We measure calibration quality by expected calibration error (ECE)~\cite{naeini2015ece} with 15 equal-width bins on the top-label confidence.

\subsection{Confidence Cascade}
\label{sec:cascade}
Calibrated probabilities enable a principled cascade: accept the CaliDec answer when the model is confident, and escalate to a more expensive model when it is not. Formally, for a decision with calibrated probability vector $\mathbf{p}$:
\begin{equation}
    \hat{y} = \begin{cases}
        \arg\max_k\; p_k & \text{if } \max_k\; p_k \ge t, \\
        \text{LLM}(q, \mathcal{O}_d, s) & \text{otherwise,}
    \end{cases}
    \label{eq:cascade}
\end{equation}
where $t$ is chosen per decision on the calibration split to maximise accuracy, with near-ties (within 0.2 points) broken in favour of fewer escalations. Because $t=0$ (never escalate) is among the candidates, the selection can decide not to use the LLM at all for a decision where it does not help. In-distribution, calibration gives $t$ a direct frequentist reading---a confidence of 0.80 means roughly 80\% of such predictions are correct.

\subsection{Baselines}
\label{sec:baselines}
We use three kinds of comparison systems. All of them are evaluated on the same test splits and receive the same post-hoc treatment as CaliDec: any score-to-option mapping and the temperature are fitted on the calibration split and applied unchanged to every test set.

\textbf{Fine-tuned 1-pass classifier (controlled ablation).} The same MiniLM-L12 backbone is fine-tuned on the same training examples, with the same epochs, learning rate, batch size and development-set checkpoint selection, but as a standard classifier: a single input \texttt{[CLS]} question-template \texttt{[SEP]} context \texttt{[SEP]}, a shared encoder and one linear head per decision type (3, 2 and 2 outputs). It needs one forward pass per decision but its options are positional, so it is neither order-invariant nor able to take unseen option sets.

\textbf{Compact groundedness checkers.} For the groundedness decision we evaluate task-specific models of comparable size: an off-the-shelf DeBERTa-v3 NLI model trained on MultiNLI, FEVER and ANLI~\cite{he2023debertav3, laurer2024less} (base, 184\,M, and large, 435\,M), AlignScore-base~\cite{zha2023alignscore} and MiniCheck (RoBERTa-large and DeBERTa-v3-large)~\cite{tang2024minicheck}. Each yields $P(\text{supported})$; long documents are split into $\sim$400-token chunks and the maximum is taken, as in the original methods. A one-feature logistic regression on the log-odds maps it to the two options. We also report their native zero-shot rule $P \ge 0.5$.

\textbf{Zero-shot references.} BGE-reranker-v2-m3 (568\,M~\cite{bge_m3}) gives one relevance score per (query or claim, context) pair, which a logistic regression maps to the decision's $K$ options. Qwen3-4B-Instruct~\cite{qwen2025qwen3} receives the question, the lettered options and the context in a chat prompt; option probabilities are read from the answer-letter logits of a single forward pass. Neither was trained on our decision formats, so they are \emph{references} that show what an off-the-shelf reranker or a prompted LLM delivers in a pipeline, not like-for-like competitors.

%═══════════════════════════════════════════════════════════════════════════════
\section{Experimental Setup}
\label{sec:experiments}
%═══════════════════════════════════════════════════════════════════════════════

\subsection{Datasets}
\label{sec:datasets}
Training examples are constructed from three public datasets, each contributing to specific decision types. Their validation sets are split into calibration and test halves. Several additional datasets, never seen during training or calibration, serve as held-out evaluation.

\begin{itemize}[leftmargin=*]
    \item \textbf{HotpotQA}~\cite{yang2018hotpotqa} (distractor setting): multi-hop questions over Wikipedia with supporting facts. Each question with two gold paragraphs and at least three distractors yields six examples: four \emph{relevance} examples (each gold paragraph containing the answer string $\to$ ``directly answers'', other gold $\to$ ``partially relevant'', two sampled distractors $\to$ ``irrelevant'') and two \emph{sufficiency} examples (both gold paragraphs plus a distractor $\to$ ``yes''; one gold paragraph swapped for a second distractor $\to$ ``no'', so the passage count is not a cue). Of 6,000 sampled training questions, 5,960 pass the filter, giving 23,840 relevance and 11,920 sufficiency examples. Calibration and test are split by question.
    \item \textbf{SQuAD~2.0}~\cite{rajpurkar2018squad2}: reading comprehension with unanswerable questions. \emph{Sufficiency} labels (answerable $\to$ ``yes''; unanswerable $\to$ ``no''), 6,000 balanced training examples.
    \item \textbf{MultiNLI}~\cite{williams2018multinli}: textual entailment. \emph{Groundedness} labels (entailment $\to$ ``yes''; contradiction or neutral $\to$ ``no''), 8,000 balanced training examples.
    \item \textbf{MuSiQue}~\cite{trivedi2022musique} (held-out): multi-hop QA with hard retrieved distractors; relevance (all hop counts) and 2-hop sufficiency, same labelling rules as HotpotQA.
    \item \textbf{VitaminC}~\cite{schuster2021vitaminc} (held-out): contrastive claim verification; groundedness, 1,000 balanced claims.
    \item \textbf{LLM-AggreFact}~\cite{tang2024minicheck} (held-out): multi-sentence, LLM-generated claims over real documents from 11 sources; groundedness, up to 100 balanced claims per source from the test split.
\end{itemize}

The training pool therefore contains $5{,}960 \times 6 + 6{,}000 + 8{,}000 = 49{,}760$ examples, from which a development split is carved (Section~\ref{sec:protocol}). Table~\ref{tab:datasplit} gives the per-decision, per-source counts of every split.

\begin{table}[H]
\caption{Examples per decision, source and split. Dev is carved from the training pool (5\%, grouped by question) and used only for checkpoint selection; Calib only for temperatures, cascade thresholds and baseline mappings; Test only for reporting. Held-out sets are never used for training or calibration.}
\label{tab:datasplit}
\centering
\small
\begin{adjustbox}{max width=\linewidth}
\begin{tabular}{llrrrr}
\toprule
\textbf{Source} & \textbf{Decision} & \textbf{Train} & \textbf{Dev} & \textbf{Calib} & \textbf{Test} \\
\midrule
HotpotQA & relevance & 22,648 & 1,192 & 1,988 & 1,984 \\
HotpotQA & sufficient & 11,324 & 596 & 994 & 992 \\
SQuAD~2.0 & sufficient & 5,700 & 300 & 500 & 500 \\
MultiNLI & grounded & 7,600 & 400 & 500 & 500 \\
\midrule
\textbf{Total} & & 47,272 & 2,488 & 3,982 & 3,976 \\
\midrule
MuSiQue (held-out) & relevance & --- & --- & --- & 4,655 \\
MuSiQue (held-out) & sufficient & --- & --- & --- & 1,040 \\
VitaminC (held-out) & grounded & --- & --- & --- & 1,000 \\
LLM-AggreFact (held-out) & grounded & --- & --- & --- & \ph{--} \\
\bottomrule
\end{tabular}

\end{adjustbox}
\end{table}

\subsection{Model Variants}
We compare the following systems:
\begin{enumerate}[leftmargin=*]
    \item \textbf{CaliDec} (proposed, revision protocol): MiniLM-L12-H384, 33.4\,M parameters, MIT licence; \NSeeds{} seeds.
    \item \textbf{1-pass classifier}: same backbone, data and protocol; \NSeeds{} seeds (Section~\ref{sec:baselines}).
    \item \textbf{CaliDec v0.1}: MiniLM-L6, 22.7\,M parameters, initialised from \texttt{cross-encoder/ms-marco-MiniLM-L6-v2} (MS MARCO pre-training).
    \item \textbf{Compact groundedness checkers}: DeBERTa-v3 NLI (base, large), AlignScore-base, MiniCheck (RoBERTa-L, DeBERTa-L).
    \item \textbf{References}: BGE-reranker-v2-m3 and prompted Qwen3-4B-Instruct.
    \item \textbf{Cascade} (CaliDec $\to$ Qwen3-4B): confidence-based routing (Eq.~\ref{eq:cascade}), all three decisions.
\end{enumerate}

Additionally, we evaluate three base-model candidates to justify the architecture choice (Section~\ref{sec:baseselection}).

\subsection{Training and Selection Protocol}
\label{sec:protocol}
CaliDec is fine-tuned from the pre-trained base for 3 epochs with AdamW~\cite{loshchilov2019decoupled} ($\beta_1=0.9$, $\beta_2=0.999$, weight decay 0.01), batch size 16 (examples, not pairs), peak learning rate $5 \times 10^{-5}$ with 6\% linear warmup and linear decay, bf16 mixed precision and gradient clipping at 1.0. \emph{Data use is strictly separated}: a development split of \NDev{} examples (5\% of the training pool, sampled by question so that all examples of one HotpotQA question fall on the same side) selects the epoch with the lowest development NLL; the remaining \NTrain{} examples are used for training; the calibration split (\NCalib{} examples) is used only after training, to fit temperatures, cascade thresholds and the baselines' score mappings; the test split (\NTest{} examples) is used only for reporting.\footnote{The submitted version selected the epoch on a slice of the calibration split. All results in this revision use the protocol above; the base-model comparison in Section~\ref{sec:baseselection} is the only exception and is labelled as such.} CaliDec and the 1-pass classifier are each trained with \NSeeds{} seeds (0--4); one run takes about 40 minutes on a single NVIDIA RTX 3060 (12\,GB).

\subsection{Evaluation Metrics and Statistics}
\begin{itemize}[leftmargin=*]
    \item \textbf{Accuracy}: fraction of decisions where the highest-probability option matches the ground truth (unchanged by temperature scaling).
    \item \textbf{Macro-F1}, \textbf{NLL}, \textbf{Brier score} and \textbf{ECE} (15 bins), after calibration unless marked ``raw''.
    \item \textbf{Latency}: milliseconds per decision, defined as the wall-clock time to score the first 480 test examples of a decision type divided by 480, after two warm-up batches, including tokenisation and excluding model loading. GPU: batches of 32 examples (CaliDec: all $K$ option pairs of the 32 examples in one forward pass; LLM: 8 prompts per forward pass, no generation), RTX 3060. CPU: batch size 1, 8 threads, Ryzen 7 5800X. The same protocol is used for every system and every table.
    \item \textbf{Statistics}: CaliDec and the 1-pass classifier are reported as mean $\pm$ standard deviation over seeds. For single systems we give bootstrap 95\% confidence intervals (10,000 resamples of test examples)~\cite{efron1993bootstrap}. Differences between two systems on the same examples are tested with the exact McNemar test~\cite{mcnemar1947note, dietterich1998approximate} and a paired bootstrap interval of the accuracy difference. With $n=\NTestGrd$ grounded test examples, the 95\% interval of a single accuracy is about $\pm$3 points, so small differences on this decision are not interpretable without the paired test.
\end{itemize}

\subsection{Hardware and Implementation}
All experiments run on a single NVIDIA GeForce RTX 3060 (12\,GB) with PyTorch 2.1+ and Transformers 4.45+. The implementation includes unit tests for: correct probability distributions from all decision types, option-order invariance, temperature fitting and calibration, command-line interface, version registry, and end-to-end network tests on all released weight versions. The code is released as the \texttt{nano-jev} Python package (PyPI), with weights on Hugging Face at \texttt{sdmlai/nano-jev}; CaliDec is the name of the method, \texttt{nano-jev} the name of its software release.

%═══════════════════════════════════════════════════════════════════════════════
\section{Results and Discussion}
\label{sec:results}
%═══════════════════════════════════════════════════════════════════════════════

\subsection{Base Model Selection}
\label{sec:baseselection}
Table~\ref{tab:baseselection} compares three clean-licence base models, all trained with the same recipe. MiniLM-L12-H384 achieves the best development NLL and test accuracy across all three decision types. Compared to the v0.1 base (MS MARCO-pretrained MiniLM-L6), the clean MiniLM-L6 is 3--5 points weaker on relevance, confirming that MS MARCO pre-training helps, but the deeper clean model more than recovers the gap. This comparison was run once under the submitted protocol (epoch selection on a calibration slice); it only motivates the choice of encoder, and no other reported number depends on it.

\begin{table}[H]
\caption{Base model selection (single run, submitted protocol). All trained with 3 epochs, lr $5\!\times\!10^{-5}$, best epoch kept. Test accuracy is after temperature calibration.}
\label{tab:baseselection}
\centering
\small
\begin{adjustbox}{max width=\linewidth}
\begin{tabular}{lcccc}
\toprule
\textbf{Base model} & \textbf{Params} & \textbf{Dev NLL} & \textbf{Test acc} & \textbf{Held-out acc} \\
 & & & (rel / suf / grd) & (rel / suf / grd) \\
\midrule
MiniLM-L6-H384 & 22.7\,M & 0.564 & .764 / .730 / .754 & .582 / .628 / .653 \\
ELECTRA-small & 13.5\,M & 0.495 & .770 / .783 / .802 & .569 / .692 / .651 \\
\textbf{MiniLM-L12-H384} & \textbf{33.4\,M} & \textbf{0.429} & \textbf{.815 / .845 / .844} & \textbf{.635} / .670 / \textbf{.675} \\
\midrule
\textit{v0.1 (MS MARCO)} & 22.7\,M & 0.491 & .798 / .759 / .804 & .631 / \textit{.680} / .672 \\
\bottomrule
\end{tabular}
\end{adjustbox}
\end{table}

\subsection{Main Results}
\label{sec:mainresults}
Table~\ref{tab:main} compares all systems on the test split and on the held-out datasets.

\begin{table*}[t]
\caption{Decision accuracy. \emph{Test}: held-out halves of HotpotQA / SQuAD~2.0 / MultiNLI validation sets. \emph{Held-out}: MuSiQue (relevance, sufficiency) and VitaminC (groundedness), never used in training or calibration. CaliDec and the 1-pass classifier: mean{\scriptsize$\pm$std} over \NSeeds{} seeds. $^\dagger$Zero-shot / adapted references (Section~\ref{sec:baselines}). Paired tests for the close comparisons are in Tables~\ref{tab:onepass} and~\ref{tab:cascadeall} and in the text.}
\label{tab:main}
\centering
\begin{adjustbox}{max width=\linewidth}
\begin{tabular}{lcccccccc}
\toprule
& & \multicolumn{3}{c}{\textbf{Test accuracy}} & \multicolumn{3}{c}{\textbf{Held-out accuracy}} \\
\cmidrule(lr){3-5}\cmidrule(lr){6-8}
\textbf{System} & \textbf{Params} & rel. & suf. & grd. & rel. & suf. & grd. \\
\midrule
BGE-reranker-v2-m3$^\dagger$ & 568\,M & \ph{--} & \ph{--} & \ph{--} & \ph{--} & \ph{--} & \ph{--} \\
Qwen3-4B-Instruct$^\dagger$ & 4\,B & \ph{--} & \ph{--} & \ph{--} & \ph{--} & \ph{--} & \ph{--} \\
\midrule
CaliDec v0.1 & 22.7\,M & 0.798 & 0.759 & 0.804 & 0.631 & 0.680 & 0.672 \\
1-pass classifier & 33.4\,M & \ph{--} & \ph{--} & \ph{--} & \ph{--} & \ph{--} & \ph{--} \\
\textbf{CaliDec (K-pass)} & 33.4\,M & \ph{--} & \ph{--} & \ph{--} & \ph{--} & \ph{--} & \ph{--} \\
\bottomrule
\end{tabular}

\end{adjustbox}
\end{table*}

\textbf{In-distribution.} CaliDec reaches \NJRelTest{} / \NJSufTest{} / \NJGroTest{} on relevance, sufficiency and groundedness, against \BGERelTest{} / \BGESufTest{} / \BGEGroTest{} for the reranker and \LLMRelTest{} / \LLMSufTest{} / \LLMGroTest{} for the prompted LLM. \revnote{Describe the gaps; for groundedness, quote the McNemar p-value of CaliDec vs.\ Qwen from results/v2/stats/nanojev\_vs\_llm\_test.md and only call it a tie/win if the test supports it.}

\textbf{Out-of-distribution.} Every model drops on unseen data. CaliDec reaches \NJRelHeld{} / \NJSufHeld{} / \NJGroHeld{} on MuSiQue relevance, MuSiQue sufficiency and VitaminC, against \LLMRelHeld{} / \LLMSufHeld{} / \LLMGroHeld{} for the LLM. \revnote{Submitted version: CaliDec led on held-out relevance and sufficiency, the LLM led on VitaminC by about 11 points; confirm with the new runs.} Contrastive claim verification appears to need more reasoning than a 33\,M-parameter encoder learns from MultiNLI; Section~\ref{sec:groundedbaselines} examines whether task-specific checkers of similar size do better.

\subsection{K-pass vs.\ 1-pass}
\label{sec:onepass}
Table~\ref{tab:onepass} isolates the K-pass design: the 1-pass classifier shares backbone, data, recipe and selection protocol with CaliDec and differs only in how options enter the model.

\begin{table*}[t]
\caption{K-pass (CaliDec) vs.\ 1-pass classifier, same backbone, data and protocol. Accuracy mean{\scriptsize$\pm$std} over seeds; $\Delta_{\text{test}}$ and McNemar $p$ compare the seed-0 models on identical test examples.}
\label{tab:onepass}
\centering
\begin{adjustbox}{max width=\linewidth}
\begin{tabular}{lcccccc}
\toprule
& \multicolumn{2}{c}{\textbf{K-pass (CaliDec)}} & \multicolumn{2}{c}{\textbf{1-pass}} & & \\
\cmidrule(lr){2-3}\cmidrule(lr){4-5}
\textbf{Decision} & Test & Held-out & Test & Held-out & $\Delta_{\text{test}}$ & McNemar $p$ \\
\midrule
relevance & \ph{--} & \ph{--} & \ph{--} & \ph{--} & \ph{--} & \ph{--} \\
sufficient & \ph{--} & \ph{--} & \ph{--} & \ph{--} & \ph{--} & \ph{--} \\
grounded & \ph{--} & \ph{--} & \ph{--} & \ph{--} & \ph{--} & \ph{--} \\
\midrule
ECE$_{\mathrm{cal}}$ (test, rel/suf/grd) & \multicolumn{2}{c}{\ph{--} / \ph{--} / \ph{--}} & \multicolumn{2}{c}{\ph{--} / \ph{--} / \ph{--}} & & \\
GPU ms/decision (rel/suf/grd) & \multicolumn{2}{c}{\ph{--} / \ph{--} / \ph{--}} & \multicolumn{2}{c}{\ph{--} / \ph{--} / \ph{--}} & & \\
Option-order invariant & \multicolumn{2}{c}{yes} & \multicolumn{2}{c}{no} & & \\
Unseen option sets & \multicolumn{2}{c}{yes (Sec.~\ref{sec:custom})} & \multicolumn{2}{c}{no (retrain)} & & \\
\bottomrule
\end{tabular}

\end{adjustbox}
\end{table*}

The 1-pass classifier needs one forward pass per decision instead of $K$ and is therefore faster (Table~\ref{tab:latency}). \revnote{Two possible readings, keep the one the data supports: (a) K-pass is significantly more accurate and/or better calibrated, so the extra passes buy accuracy; (b) the two are within noise, in which case the case for K-pass rests on order invariance and a single model for all option sets, and accuracy is not claimed as an advantage of the design.}

\subsection{Compact Groundedness Checkers}
\label{sec:groundedbaselines}
Table~\ref{tab:groundedbaselines} compares CaliDec with task-specific groundedness models on the in-distribution test split (MultiNLI), on VitaminC, and on LLM-AggreFact, whose claims are multi-sentence LLM outputs over long documents and thus closer to RAG answer checking.

\begin{table}[H]
\caption{Groundedness accuracy of CaliDec and compact task-specific checkers. All mappings and temperatures fitted on the MultiNLI calibration split and applied unchanged. ``chunk-max'': documents split into $\sim$400-token chunks, maximum $P(\text{supported})$ over chunks. $^\dagger$References.}
\label{tab:groundedbaselines}
\centering
\small
\begin{adjustbox}{max width=\linewidth}
\begin{tabular}{lcccc}
\toprule
\textbf{System} & \textbf{Params} & \textbf{MultiNLI (test)} & \textbf{VitaminC} & \textbf{LLM-AggreFact} \\
\midrule
CaliDec (K-pass, trunc. 512) & 33\,M & \ph{--} & \ph{--} & \ph{--} \\
CaliDec (K-pass, chunk-max) & 33\,M & \ph{--} & \ph{--} & \ph{--} \\
1-pass classifier & 33\,M & \ph{--} & \ph{--} & \ph{--} \\
DeBERTa-v3-base NLI & 184\,M & \ph{--} & \ph{--} & \ph{--} \\
DeBERTa-v3-large NLI & 435\,M & \ph{--} & \ph{--} & \ph{--} \\
AlignScore-base & 125\,M & \ph{--} & \ph{--} & \ph{--} \\
MiniCheck-RoBERTa-L & 355\,M & \ph{--} & \ph{--} & \ph{--} \\
MiniCheck-DeBERTa-L & 435\,M & \ph{--} & \ph{--} & \ph{--} \\
BGE-reranker-v2-m3$^\dagger$ & 568\,M & \ph{--} & \ph{--} & \ph{--} \\
Qwen3-4B-Instruct$^\dagger$ & 4\,B & \ph{--} & \ph{--} & \ph{--} \\
\bottomrule
\end{tabular}

\end{adjustbox}
\end{table}

\revnote{Interpret: MiniCheck and the NLI models were trained on far more (and more varied) entailment data, so they are expected to lead on VitaminC and AggreFact. If CaliDec trails them, say so plainly and position CaliDec's groundedness head as the calibrated, cascade-ready member of a multi-decision model rather than the best stand-alone checker.}

\subsection{Calibration Quality}
\label{sec:calresults}
Table~\ref{tab:calibration} reports ECE before and after temperature scaling.

\begin{table}[H]
\caption{Calibration. ECE before (raw) and after temperature scaling (cal); temperatures fitted on the calibration split and applied unchanged to the held-out sets. $\Delta$ECE$_{\mathrm{cal}}$: held-out minus test. CaliDec and 1-pass: mean over seeds.}
\label{tab:calibration}
\centering
\small
\begin{adjustbox}{max width=\linewidth}
\begin{tabular}{llcccccc}
\toprule
& & \multicolumn{3}{c}{\textbf{Test}} & \multicolumn{3}{c}{\textbf{Held-out}} \\
\cmidrule(lr){3-5}\cmidrule(lr){6-8}
\textbf{System} & \textbf{Dec.} & ECE$_{\mathrm{raw}}$ & ECE$_{\mathrm{cal}}$ & $\tau$ & ECE$_{\mathrm{raw}}$ & ECE$_{\mathrm{cal}}$ & $\Delta$ECE$_{\mathrm{cal}}$ \\
\midrule
\multirow{3}{*}{CaliDec} & rel. & \ph{--} & \ph{--} & \ph{--} & \ph{--} & \ph{--} & \ph{--} \\
 & suf. & \ph{--} & \ph{--} & \ph{--} & \ph{--} & \ph{--} & \ph{--} \\
 & grd. & \ph{--} & \ph{--} & \ph{--} & \ph{--} & \ph{--} & \ph{--} \\
\midrule
\multirow{3}{*}{1-pass} & rel. & \ph{--} & \ph{--} & \ph{--} & \ph{--} & \ph{--} & \ph{--} \\
 & suf. & \ph{--} & \ph{--} & \ph{--} & \ph{--} & \ph{--} & \ph{--} \\
 & grd. & \ph{--} & \ph{--} & \ph{--} & \ph{--} & \ph{--} & \ph{--} \\
\midrule
\multirow{3}{*}{CaliDec v0.1} & rel. & 0.061 & 0.020 & 1.45 & 0.223 & 0.163 & +0.143 \\
 & suf. & 0.065 & 0.019 & 1.38 & 0.126 & 0.081 & +0.062 \\
 & grd. & 0.086 & 0.045 & 1.23 & 0.202 & 0.177 & +0.132 \\
\midrule
\multirow{3}{*}{BGE-reranker} & rel. & \ph{--} & \ph{--} & \ph{--} & \ph{--} & \ph{--} & \ph{--} \\
 & suf. & \ph{--} & \ph{--} & \ph{--} & \ph{--} & \ph{--} & \ph{--} \\
 & grd. & \ph{--} & \ph{--} & \ph{--} & \ph{--} & \ph{--} & \ph{--} \\
\midrule
\multirow{3}{*}{Qwen3-4B} & rel. & \ph{--} & \ph{--} & \ph{--} & \ph{--} & \ph{--} & \ph{--} \\
 & suf. & \ph{--} & \ph{--} & \ph{--} & \ph{--} & \ph{--} & \ph{--} \\
 & grd. & \ph{--} & \ph{--} & \ph{--} & \ph{--} & \ph{--} & \ph{--} \\
\bottomrule
\end{tabular}

\end{adjustbox}
\end{table}

In-distribution, temperature scaling brings CaliDec's ECE to \NJEceTestMin--\NJEceTestMax. On the held-out sets the same temperatures give ECE \NJEceHeldMin--\NJEceHeldMax: the probabilities overstate accuracy on unseen domains, as expected from a single scalar that cannot correct a distribution shift~\cite{guo2017calibration}. This is the main practical caveat of the paper, and it has a direct consequence for deployment: temperatures (and cascade thresholds) should be re-fitted on a few hundred labelled examples from the target domain. \revnote{Optional evidence: re-fitting on the LLM-AggreFact dev split (results/v2/nanojev/seed0/aggrefact\_refit.json) vs.\ the in-domain temperature; one sentence with the two ECE values.}

\subsection{Confidence Cascade}
\label{sec:cascaderesults}
Table~\ref{tab:cascade} gives the full threshold sweep for groundedness, the decision on which the LLM is strongest; Table~\ref{tab:cascadeall} gives the cascade at the calibration-selected threshold for all three decisions.

\begin{table}[H]
\caption{Confidence cascade CaliDec $\to$ Qwen3-4B on the \emph{grounded} decision (seed-0 model). \% esc.: fraction of decisions escalated to the LLM, reported separately for each test set. The chosen threshold is selected on the calibration split. ms/dec: CaliDec latency plus the escalated fraction times the LLM latency.}
\label{tab:cascade}
\centering
\small
\begin{adjustbox}{max width=\linewidth}
\begin{tabular}{lccccc}
\toprule
& \multicolumn{3}{c}{\textbf{Test (MultiNLI)}} & \multicolumn{2}{c}{\textbf{Held-out (VitaminC)}} \\
\cmidrule(lr){2-4}\cmidrule(lr){5-6}
\textbf{Strategy} & \textbf{\% esc.} & \textbf{Acc} & \textbf{ms/dec} & \textbf{\% esc.} & \textbf{Acc} \\
\midrule
CaliDec only & \ph{--} & \ph{--} & \ph{--} & \ph{--} & \ph{--} \\
Cascade, $t=0.60$ & \ph{--} & \ph{--} & \ph{--} & \ph{--} & \ph{--} \\
Cascade, $t=0.70$ & \ph{--} & \ph{--} & \ph{--} & \ph{--} & \ph{--} \\
Cascade, $t=0.80$ & \ph{--} & \ph{--} & \ph{--} & \ph{--} & \ph{--} \\
Cascade, $t=0.85$ & \ph{--} & \ph{--} & \ph{--} & \ph{--} & \ph{--} \\
Cascade, $t=0.90$ & \ph{--} & \ph{--} & \ph{--} & \ph{--} & \ph{--} \\
Cascade, $t=0.95$ & \ph{--} & \ph{--} & \ph{--} & \ph{--} & \ph{--} \\
LLM only & \ph{--} & \ph{--} & \ph{--} & \ph{--} & \ph{--} \\
\bottomrule
\end{tabular}

\end{adjustbox}
\end{table}

\begin{table*}[t]
\caption{Cascade on all three decisions at the threshold $t$ chosen on the calibration split ($t=0$: never escalate). Exact McNemar $p$-values compare the cascade with each model alone on the same examples.}
\label{tab:cascadeall}
\centering
\begin{adjustbox}{max width=\linewidth}
\begin{tabular}{lccccccc}
\toprule
\textbf{Decision / set} & $t$ & \textbf{\% esc.} & \textbf{Cascade} & \textbf{CaliDec} & \textbf{LLM} & $p_{\text{vs CaliDec}}$ & $p_{\text{vs LLM}}$ \\
\midrule
relevance / test & \ph{--} & \ph{--} & \ph{--} & \ph{--} & \ph{--} & \ph{--} & \ph{--} \\
relevance / held-out & \ph{--} & \ph{--} & \ph{--} & \ph{--} & \ph{--} & \ph{--} & \ph{--} \\
sufficient / test & \ph{--} & \ph{--} & \ph{--} & \ph{--} & \ph{--} & \ph{--} & \ph{--} \\
sufficient / held-out & \ph{--} & \ph{--} & \ph{--} & \ph{--} & \ph{--} & \ph{--} & \ph{--} \\
grounded / test & \ph{--} & \ph{--} & \ph{--} & \ph{--} & \ph{--} & \ph{--} & \ph{--} \\
grounded / held-out & \ph{--} & \ph{--} & \ph{--} & \ph{--} & \ph{--} & \ph{--} & \ph{--} \\
\bottomrule
\end{tabular}

\end{adjustbox}
\end{table*}

\textbf{Groundedness, in-distribution.} At the selected threshold $t=\CascT$ the cascade reaches \CascGrdTestAcc{} with \CascGrdTestEsc{} of decisions escalated, at \CascGrdMs{}\,ms per decision versus \CascGrdLLMMs{}\,ms for the LLM alone (\CascGrdSpeedup$\times$). \revnote{Only claim that the cascade ``surpasses both models'' if both McNemar p-values in Table~\ref{tab:cascadeall} are below 0.05; otherwise say it matches the better model at a fraction of the cost.}

\textbf{Groundedness, held-out.} On VitaminC the same threshold escalates \CascGrdHeldEsc{} of decisions and reaches \CascGrdHeldAcc{}, against \CascGrdNanoHeldAcc{} for CaliDec alone and \CascGrdLLMHeldAcc{} for the LLM alone. \revnote{If the cascade is below the LLM here (submitted version: 0.745 vs.\ 0.787), state it explicitly:} Out of distribution the cascade trades accuracy for latency relative to the LLM, and reaches the LLM's accuracy only from $t=\CascGrdHeldParityT$, i.e.\ with a much larger escalated fraction. Because CaliDec's held-out probabilities are overconfident (Section~\ref{sec:calresults}), a threshold chosen in-distribution escalates too little on a new domain---one more reason to re-fit on target-domain data.

\textbf{Relevance and sufficiency.} Here the prompted LLM is clearly weaker than CaliDec (Table~\ref{tab:main}), and the threshold selected on the calibration split is \revnote{expected: $t=0$ or close to it}, so the cascade escalates (almost) nothing and its accuracy equals CaliDec's. We regard this as a desirable property rather than a null result: because thresholds are selected on calibrated probabilities with ``never escalate'' as a candidate, the cascade does not pay LLM latency to replace good answers with worse ones.

\textbf{Is calibration load-bearing?} Table~\ref{tab:ablationtemp} repeats the grounded cascade with uncalibrated CaliDec probabilities ($\tau=1$), both at a fixed $t=0.80$ and with the threshold re-selected on the calibration split.

\begin{table}[H]
\caption{Temperature ablation for the grounded cascade (test split, seed-0 model). ``raw'': $\tau=1$; ``cal.'': fitted temperature.}
\label{tab:ablationtemp}
\centering
\small
\begin{adjustbox}{max width=\linewidth}
\begin{tabular}{lcccc}
\toprule
\textbf{Strategy} & $\tau$ & $t$ & \textbf{Acc} & \textbf{\% escalated} \\
\midrule
Cascade, raw, $t=0.80$ & 1.00 & 0.80 & \ph{--} & \ph{--} \\
Cascade, cal., $t=0.80$ & \ph{--} & 0.80 & \ph{--} & \ph{--} \\
Cascade, raw, $t$ chosen on calib & 1.00 & \ph{--} & \ph{--} & \ph{--} \\
Cascade, cal., $t$ chosen on calib & \ph{--} & \ph{--} & \ph{--} & \ph{--} \\
\bottomrule
\end{tabular}

\end{adjustbox}
\end{table}

\revnote{Interpret. At a fixed threshold, the overconfident raw model escalates fewer decisions. If re-selecting $t$ on calib closes the gap, say that calibration mainly makes a \emph{fixed, interpretable} threshold work, rather than claiming it is indispensable.}

\subsection{Per-Decision Detailed Metrics}
\label{sec:detailed}
Table~\ref{tab:detailed} gives the full metric suite for CaliDec, and Table~\ref{tab:perclass} breaks the 3-way relevance decision down by class. Both tables and Table~\ref{tab:main} come from the same evaluation runs. ``Partially relevant''---a gold paragraph that does not itself contain the answer---is the hardest class, consistent with the ambiguity of the HotpotQA-derived label.

\begin{table}[H]
\caption{Detailed per-decision metrics for CaliDec on the test split, mean{\scriptsize$\pm$std} over seeds, after calibration. $\tau$: fitted temperature (mean).}
\label{tab:detailed}
\centering
\small
\begin{adjustbox}{max width=\linewidth}
\begin{tabular}{lcccccc}
\toprule
\textbf{Decision} & $n$ & $\tau$ & \textbf{Acc} & \textbf{Macro-F1} & \textbf{NLL} & \textbf{Brier} \\
\midrule
relevance & \ph{--} & \ph{--} & \ph{--} & \ph{--} & \ph{--} & \ph{--} \\
sufficient & \ph{--} & \ph{--} & \ph{--} & \ph{--} & \ph{--} & \ph{--} \\
grounded & \ph{--} & \ph{--} & \ph{--} & \ph{--} & \ph{--} & \ph{--} \\
\bottomrule
\end{tabular}

\end{adjustbox}
\end{table}

\begin{table}[H]
\caption{Per-class precision, recall and F1 for the 3-way relevance decision (CaliDec, test split, mean{\scriptsize$\pm$std} over seeds). The macro average equals the relevance macro-F1 in Table~\ref{tab:detailed}.}
\label{tab:perclass}
\centering
\small
\begin{adjustbox}{max width=\linewidth}
\begin{tabular}{lcccc}
\toprule
\textbf{Class} & \textbf{Precision} & \textbf{Recall} & \textbf{F1} & \textbf{Support} \\
\midrule
irrelevant & \ph{--} & \ph{--} & \ph{--} & \ph{--} \\
partially relevant & \ph{--} & \ph{--} & \ph{--} & \ph{--} \\
directly answers & \ph{--} & \ph{--} & \ph{--} & \ph{--} \\
\midrule
\textbf{Macro avg} & \ph{--} & \ph{--} & \ph{--} & \\
\bottomrule
\end{tabular}

\end{adjustbox}
\end{table}

\subsection{Latency}
\label{sec:latency}
Table~\ref{tab:latency} reports latency per decision type under the single protocol of Section~\ref{sec:experiments}.

\begin{table}[H]
\caption{Latency in ms per decision, per decision type. GPU: RTX 3060, batches of 32 examples (LLM: 8). CPU: Ryzen 7 5800X, 8 threads, batch size 1. Same 480 examples per decision for every system.}
\label{tab:latency}
\centering
\small
\begin{adjustbox}{max width=\linewidth}
\begin{tabular}{lccccccc}
\toprule
& & \multicolumn{3}{c}{\textbf{GPU ms/decision}} & \multicolumn{3}{c}{\textbf{CPU ms/decision}} \\
\cmidrule(lr){3-5}\cmidrule(lr){6-8}
\textbf{System} & \textbf{Params} & rel. & suf. & grd. & rel. & suf. & grd. \\
\midrule
CaliDec (K-pass) & \ph{--} & \ph{--} & \ph{--} & \ph{--} & \ph{--} & \ph{--} & \ph{--} \\
1-pass classifier & \ph{--} & \ph{--} & \ph{--} & \ph{--} & \ph{--} & \ph{--} & \ph{--} \\
BGE-reranker-v2-m3 & \ph{--} & \ph{--} & \ph{--} & \ph{--} & \ph{--} & \ph{--} & \ph{--} \\
Qwen3-4B-Instruct & \ph{--} & \ph{--} & \ph{--} & \ph{--} & \ph{--} & \ph{--} & \ph{--} \\
\bottomrule
\end{tabular}

\end{adjustbox}
\end{table}

Latency differs by decision type mainly through the number of options $K$ (3 for relevance, 2 otherwise) and the context length (sufficiency contexts hold three passages). On GPU, CaliDec is \SpeedupMin--\SpeedupMax$\times$ faster than the prompted LLM, the ratio being taken per decision type. \revnote{Add one sentence on 1-pass vs.\ K-pass latency.}

\subsection{Effect of the 512-Token Limit}
\label{sec:truncation}
Table~\ref{tab:truncation} reports how often inputs exceed 512 tokens, how often a gold passage is cut, and accuracy on truncated vs.\ non-truncated inputs. In-distribution, at most \TruncInDistMax{} of inputs of any decision are truncated, so the limit does not affect the main results. The exception is MuSiQue sufficiency, where \TruncMusSufPct{} of inputs are truncated and \TruncMusSufGold{} lose part of a gold passage; accuracy is \TruncMusSufAccT{} on truncated vs.\ \TruncMusSufAccNT{} on non-truncated inputs. Truncating each passage to an equal share of the budget, so that every passage survives, gives \TruncMusSufAccBal{} on MuSiQue sufficiency. \revnote{Interpret the balanced variant.}

\begin{table*}[t]
\caption{Truncation at 512 tokens (seed-0 model). \% gold cut: multi-passage states in which a gold passage (labelled relevant for the same query) is partly or fully cut. Acc trunc.\ (count) / not trunc.: accuracy on the two subsets. Acc balanced: accuracy when each passage is cut to an equal share of the budget.}
\label{tab:truncation}
\centering
\begin{adjustbox}{max width=\linewidth}
\begin{tabular}{lcccccc}
\toprule
\textbf{Decision / source} & $n$ & \textbf{\% $>$512} & \textbf{\% gold cut} & \textbf{Acc trunc.} & \textbf{Acc not trunc.} & \textbf{Acc balanced} \\
\midrule
\ph{--} & & & & & & \\
\bottomrule
\end{tabular}

\end{adjustbox}
\end{table*}

\subsection{End-to-End RAG Evaluation}
\label{sec:e2e}
To measure downstream effect, we run a small RAG pipeline on 300 HotpotQA validation questions and 300 MuSiQue questions that appear in none of our calibration, test or held-out files. For each question, BM25~\cite{robertson2009bm25} ranks the question's own paragraph pool (10 paragraphs for HotpotQA, 20 for MuSiQue) and keeps the top 5; a filter then removes passages judged irrelevant---no filter, the BGE reranker with the calibration-fitted mapping, CaliDec's relevance decision, or the prompted Qwen3-4B relevance judge---and Qwen3-4B-Instruct generates a short answer from the remaining passages (greedy decoding; if a filter removes everything, the BM25 top-1 is kept). Table~\ref{tab:e2e} reports exact match and F1, recall and precision of gold passages in the kept set, and total latency per question (retrieval + filter + generation).

\begin{table}[H]
\caption{End-to-end RAG on 300 HotpotQA and 300 MuSiQue questions (seed-0 CaliDec). Gold rec./prec.: recall and precision of the gold paragraphs among the passages passed to the generator. ms/q: retrieval + filter + generation per question.}
\label{tab:e2e}
\centering
\small
\begin{adjustbox}{max width=\linewidth}
\begin{tabular}{llcccccc}
\toprule
\textbf{Set} & \textbf{Filter} & \textbf{EM} & \textbf{F1} & \textbf{Gold rec.} & \textbf{Gold prec.} & \textbf{Kept} & \textbf{ms/q} \\
\midrule
HotpotQA & none (top-5) & \ph{--} & \ph{--} & \ph{--} & \ph{--} & \ph{--} & \ph{--} \\
HotpotQA & BGE + LR & \ph{--} & \ph{--} & \ph{--} & \ph{--} & \ph{--} & \ph{--} \\
HotpotQA & CaliDec & \ph{--} & \ph{--} & \ph{--} & \ph{--} & \ph{--} & \ph{--} \\
HotpotQA & Qwen3-4B judge & \ph{--} & \ph{--} & \ph{--} & \ph{--} & \ph{--} & \ph{--} \\
\midrule
MuSiQue & none (top-5) & \ph{--} & \ph{--} & \ph{--} & \ph{--} & \ph{--} & \ph{--} \\
MuSiQue & BGE + LR & \ph{--} & \ph{--} & \ph{--} & \ph{--} & \ph{--} & \ph{--} \\
MuSiQue & CaliDec & \ph{--} & \ph{--} & \ph{--} & \ph{--} & \ph{--} & \ph{--} \\
MuSiQue & Qwen3-4B judge & \ph{--} & \ph{--} & \ph{--} & \ph{--} & \ph{--} & \ph{--} \\
\bottomrule
\end{tabular}

\end{adjustbox}
\end{table}

With the CaliDec filter, F1 is \EEHotNJFone{} on HotpotQA and \EEMusNJFone{} on MuSiQue, against \EEHotNoneFone{} / \EEMusNoneFone{} without a filter and \EEHotQwenFone{} / \EEMusQwenFone{} with the LLM judge; total latency per question is \EEHotNJMs{} vs.\ \EEHotQwenMs{}\,ms on HotpotQA. \revnote{Interpret: does filtering help the generator, and is CaliDec close to the LLM judge at a fraction of the filter latency? This is a small study; say so.}

\subsection{Custom Option Sets}
\label{sec:custom}
Because options are scored independently, any option set can be passed to the model. To test whether this works without retraining, we evaluate zero-shot topic and emotion classification with label sets never seen in training: AG News ($K=4$), Yahoo Answers topics ($K=10$)~\cite{zhang2015character} and Emotion ($K=6$)~\cite{saravia2018carer}, 1,000 test texts each (889 for Emotion), with the label names as options.

\begin{table}[H]
\caption{Zero-shot accuracy on unseen option sets. Random: $1/K$. The probe row is a \emph{supervised} reference: a logistic regression on frozen 1-pass-encoder features trained on 300 labelled examples per task.}
\label{tab:custom}
\centering
\small
\begin{adjustbox}{max width=\linewidth}
\begin{tabular}{lcccccc}
\toprule
\textbf{Task} & $K$ & \textbf{Random} & \textbf{v0.1} & \textbf{v1.0} & \textbf{CaliDec (rev.)} & \textbf{Qwen3-4B} \\
\midrule
AG News & 4 & 0.250 & 0.396 & 0.186 & \ph{--} & \ph{--} \\
Yahoo Answers & 10 & 0.100 & 0.239 & 0.056 & \ph{--} & \ph{--} \\
Emotion & 6 & 0.167 & 0.334 & 0.138 & \ph{--} & \ph{--} \\
\midrule
\multicolumn{7}{l}{\textit{Supervised reference: linear probe on frozen 1-pass encoder, 300 labels}} \\
AG News (probe) & 4 & & & & \ph{--} &  \\
Yahoo Answers (probe) & 10 & & & & \ph{--} &  \\
Emotion (probe) & 6 & & & & \ph{--} &  \\
\bottomrule
\end{tabular}

\end{adjustbox}
\end{table}

The results do not support custom option sets as a capability of the released models. v0.1, whose base was pre-trained on MS MARCO, is above chance on all three tasks (\CustAgnewsVZero{} vs.\ \CustAgnewsRand{} on AG News, \CustYahooVZero{} vs.\ \CustYahooRand{} on Yahoo, \CustEmotionVZero{} vs.\ \CustEmotionRand{} on Emotion) but far from useful, and v1.0 is \emph{below} chance on all three (\CustAgnewsVOne{}, \CustYahooVOne{}, \CustEmotionVOne{}): trained only on three fixed option vocabularies, the clean-licence model scores options by surface familiarity rather than by meaning. We therefore claim order invariance and a uniform interface for the three trained decisions, and leave option-set generalisation---for example training with paraphrased option names and generic text--label matching data---to future work. \revnote{If you do the option-augmentation retrain, report it in the CaliDec (rev.) column and revise this paragraph.}

\subsection{Ablation: Training Data Size}
\label{sec:ablation_data}
\begin{table}[H]
\caption{Effect of training data size on CaliDec test accuracy (seed 0, random fraction of the training split; dev, calib and test unchanged).}
\label{tab:ablation_data}
\centering
\small
\begin{adjustbox}{max width=\linewidth}
\begin{tabular}{lccc}
\toprule
\textbf{Training fraction} & \textbf{Rel. acc} & \textbf{Suf. acc} & \textbf{Grd. acc} \\
\midrule
10\% & \ph{--} & \ph{--} & \ph{--} \\
25\% & \ph{--} & \ph{--} & \ph{--} \\
50\% & \ph{--} & \ph{--} & \ph{--} \\
100\% (seed 0) & \ph{--} & \ph{--} & \ph{--} \\
\bottomrule
\end{tabular}

\end{adjustbox}
\end{table}

\revnote{Interpret Table~\ref{tab:ablation_data}; compare the smallest fraction with the reference baselines only if the numbers support it.}

\subsection{Discussion}
\label{sec:discussion}

\textbf{Why typed decisions?} By fixing the question template and option set per decision type, we eliminate the ambiguity of free-form self-assessment and the parsing fragility of LLM-generated judgments. The model cannot hallucinate a novel option or produce a malformed answer. The schema also enables systematic evaluation: each decision type has a clear accuracy metric over its options.

\textbf{Why calibration matters.} In a RAG pipeline, decision probabilities drive downstream logic: a relevance probability determines whether a passage enters the prompt, a sufficiency probability whether to retrieve more, and a groundedness probability whether to accept or escalate. Calibration makes these thresholds interpretable in-distribution and lets the cascade threshold be chosen once on held-out data. Its limits under domain shift (Section~\ref{sec:calresults}) are the main deployment caveat.

\textbf{Limitations.}
\begin{itemize}[leftmargin=*]
    \item English, Wikipedia-style text only for relevance and sufficiency; code, tables and non-English text are untested.
    \item Temperatures and thresholds fitted in-distribution do not transfer; a few hundred target-domain labels are needed for re-fitting.
    \item Zero-shot custom option sets are unreliable in the released models (Section~\ref{sec:custom}).
    \item The 512-token limit truncates a noticeable share of multi-passage MuSiQue inputs (Section~\ref{sec:truncation}); long-document RAG would need a chunked or hierarchical design.
    \item The end-to-end evaluation is small (600 questions, one generator) and uses the datasets' own paragraph pools rather than open-corpus retrieval.
    \item The ``partially relevant'' label follows HotpotQA conventions, which may not match all use cases.
    \item All experiments use a single consumer GPU; behaviour under production-scale batching and quantisation is not measured.
\end{itemize}

%═══════════════════════════════════════════════════════════════════════════════
\section{Conclusion}
\label{sec:conclusion}
%═══════════════════════════════════════════════════════════════════════════════

We introduced CaliDec, a 33.4\,M-parameter calibrated cross-encoder that casts the recurring micro-decisions of RAG pipelines as typed, closed-set classification problems. By scoring (question, option, context) triples independently and applying per-decision temperature scaling under a strict train / dev / calibration / test separation, CaliDec produces calibrated probabilities that make thresholds interpretable and support confidence cascading to a larger model.

On in-distribution data, CaliDec reaches \NJRelTest{} / \NJSufTest{} / \NJGroTest{} accuracy on relevance, sufficiency and groundedness at \SpeedupMin--\SpeedupMax$\times$ lower latency than a prompted 4\,B-parameter LLM. \revnote{One sentence on the 1-pass comparison and one on the compact groundedness checkers, as the final numbers dictate.} With a threshold chosen on the calibration split, the cascade escalates only where the LLM helps. On unseen domains accuracy drops and calibration degrades, so temperatures and thresholds should be re-fitted on target-domain data before deployment.

Future work includes: (i)~training with varied option vocabularies so that custom \texttt{decide()} calls generalise; (ii)~a \texttt{next\_action} decision (answer / retrieve more / rewrite / decompose / abstain) for agentic RAG loops; (iii)~a model that reads the query and all chunks in one pass, removing the 512-token truncation; and (iv)~domain-adaptive calibration~\cite{zhao2021calibrate} that transfers temperatures with few labelled examples.

%═══════════════════════════════════════════════════════════════════════════════
% Back matter
%═══════════════════════════════════════════════════════════════════════════════

\par
\vspace{10pt}
{\large\textbf{Data Availability Statement}}

The complete CaliDec implementation, including model definitions, training and evaluation scripts, baseline code, calibration routines, statistical tests and unit tests, is publicly available at \url{https://github.com/shubham10divakar/nano-jev}. Pre-trained weights are hosted on Hugging Face at \url{https://huggingface.co/sdmlai/nano-jev}. The training datasets (HotpotQA, SQuAD~2.0, MultiNLI) and evaluation datasets (MuSiQue, VitaminC, LLM-AggreFact, AG News, Yahoo Answers, Emotion) are publicly available from their original providers under their respective licences.

\acknowledgement
This work was supported without any funding.

\conflictsofinterest
The authors declare no conflicts of interest.

\vspace{10pt}

{\large\textbf{Ethical Approval and Consent to Participate}}

Not applicable.

\begin{fullwidth}
\begin{thebibliography}{99}

\bibitem{lewis2020rag}
P. Lewis, E. Perez, A. Piktus, F. Petroni, V. Karpukhin, N. Goyal, \textit{et al.}, ``Retrieval-augmented generation for knowledge-intensive NLP tasks,'' in \textit{Proc. NeurIPS}, 2020, pp. 9459--9474.

\bibitem{guu2020realm}
K. Guu, K. Lee, Z. Tung, P. Pasupat, and M.-W. Chang, ``Retrieval augmented language model pre-training,'' in \textit{Proc. ICML}, 2020, pp. 3929--3938.

\bibitem{asai2024selfrag}
A. Asai, Z. Wu, Y. Wang, A. Sil, and H. Hajishirzi, ``Self-RAG: Learning to retrieve, generate, and critique through self-reflection,'' in \textit{Proc. ICLR}, 2024.

\bibitem{yan2024corrective}
S. Yan, S. Gu, L. Zhu, and Z. Dai, ``Corrective retrieval augmented generation,'' \textit{arXiv:2401.15884}, 2024.

\bibitem{kadavath2022language}
S. Kadavath, T. Conerly, A. Askell, T. Henighan, D. Drain, E. Perez, \textit{et al.}, ``Language models (mostly) know what they know,'' \textit{arXiv:2207.05221}, 2022.

\bibitem{kuhn2023semantic}
L. Kuhn, Y. Gal, and S. Farquhar, ``Semantic uncertainty: Linguistic invariances for uncertainty estimation in natural language generation,'' in \textit{Proc. ICLR}, 2023.

\bibitem{bge_m3}
J. Chen, S. Xiao, P. Zhang, K. Luo, D. Lian, and Z. Liu, ``BGE M3-Embedding: Multi-lingual, multi-functionality, multi-granularity text embeddings through self-knowledge distillation,'' \textit{arXiv:2402.03216}, 2024.

\bibitem{nogueira2019passage}
R. Nogueira, W. Yang, K. Cho, and J. Lin, ``Multi-stage document ranking with BERT,'' \textit{arXiv:1910.14424}, 2019.

\bibitem{reimers2019sentence}
N. Reimers and I. Gurevych, ``Sentence-BERT: Sentence embeddings using Siamese BERT-networks,'' in \textit{Proc. EMNLP-IJCNLP}, 2019, pp. 3982--3992.

\bibitem{devlin2019bert}
J. Devlin, M.-W. Chang, K. Lee, and K. Toutanova, ``BERT: Pre-training of deep bidirectional transformers for language understanding,'' in \textit{Proc. NAACL-HLT}, 2019, pp. 4171--4186.

\bibitem{liu2019roberta}
Y. Liu, M. Ott, N. Goyal, J. Du, M. Joshi, D. Chen, O. Levy, M. Lewis, L. Zettlemoyer, and V. Stoyanov, ``RoBERTa: A robustly optimized BERT pretraining approach,'' \textit{arXiv:1907.11692}, 2019.

\bibitem{he2023debertav3}
P. He, J. Gao, and W. Chen, ``DeBERTaV3: Improving DeBERTa using ELECTRA-style pre-training with gradient-disentangled embedding sharing,'' in \textit{Proc. ICLR}, 2023.

\bibitem{laurer2024less}
M. Laurer, W. van Atteveldt, A. Casas, and K. Welbers, ``Less annotating, more classifying: Addressing the data scarcity issue of supervised machine learning with deep transfer learning and BERT-NLI,'' \textit{Political Analysis}, vol. 32, no. 1, pp. 84--100, 2024.

\bibitem{laban2022summac}
P. Laban, T. Schnabel, P. N. Bennett, and M. A. Hearst, ``SummaC: Re-visiting NLI-based models for inconsistency detection in summarization,'' \textit{Transactions of the ACL}, vol. 10, pp. 163--177, 2022.

\bibitem{honovich2022true}
O. Honovich, R. Aharoni, J. Herzig, H. Taitelbaum, D. Kukliansy, V. Cohen, T. Scialom, I. Szpektor, A. Hassidim, and Y. Matias, ``TRUE: Re-evaluating factual consistency evaluation,'' in \textit{Proc. NAACL-HLT}, 2022, pp. 3905--3920.

\bibitem{zha2023alignscore}
Y. Zha, Y. Yang, R. Li, and Z. Hu, ``AlignScore: Evaluating factual consistency with a unified alignment function,'' in \textit{Proc. ACL}, 2023, pp. 11328--11348.

\bibitem{tang2024minicheck}
L. Tang, P. Laban, and G. Durrett, ``MiniCheck: Efficient fact-checking of LLMs on grounding documents,'' in \textit{Proc. EMNLP}, 2024.

\bibitem{es2024ragas}
S. Es, J. James, L. Espinosa-Anke, and S. Schockaert, ``RAGAS: Automated evaluation of retrieval augmented generation,'' in \textit{Proc. EACL (System Demonstrations)}, 2024, pp. 150--158.

\bibitem{saadfalcon2024ares}
J. Saad-Falcon, O. Khattab, C. Potts, and M. Zaharia, ``ARES: An automated evaluation framework for retrieval-augmented generation systems,'' in \textit{Proc. NAACL-HLT}, 2024.

\bibitem{guo2017calibration}
C. Guo, G. Pleiss, Y. Sun, and K. Q. Weinberger, ``On calibration of modern neural networks,'' in \textit{Proc. ICML}, 2017, pp. 1321--1330.

\bibitem{platt1999probabilistic}
J. C. Platt, ``Probabilistic outputs for support vector machines and comparisons to regularized likelihood methods,'' in \textit{Advances in Large Margin Classifiers}, MIT Press, 2000, pp. 61--74.

\bibitem{chen2023frugalgpt}
L. Chen, M. Zaharia, and J. Zou, ``FrugalGPT: How to use large language models while reducing cost and improving performance,'' \textit{arXiv:2305.05176}, 2023.

\bibitem{ding2024hybrid}
D. Ding, A. Mallick, C. Wang, R. Sim, S. Mukherjee, V. Raman, A. Awadallah, and C. Gao, ``Hybrid LLM: Cost-efficient and quality-aware query routing,'' in \textit{Proc. ICLR}, 2024.

\bibitem{leviathan2023fast}
Y. Leviathan, M. Kalman, and Y. Matias, ``Fast inference from transformers via speculative decoding,'' in \textit{Proc. ICML}, 2023, pp. 19274--19286.

\bibitem{jpalmer2026jev}
J. Palmer, ``Jev: A System~1 decision model for agentic AI,'' TypeSafe AI, 2026. [Online]. Available: \url{https://typesafe.ai/jev}

\bibitem{wang2020minilm}
W. Wang, F. Wei, L. Dong, H. Bao, N. Yang, and M. Zhou, ``MiniLM: Deep self-attention distillation for task-agnostic compression of pre-trained transformers,'' in \textit{Proc. NeurIPS}, 2020, pp. 5776--5788.

\bibitem{yang2018hotpotqa}
Z. Yang, P. Qi, S. Zhang, Y. Bengio, W. W. Cohen, R. Salakhutdinov, and C. D. Manning, ``HotpotQA: A dataset for diverse, explainable multi-hop question answering,'' in \textit{Proc. EMNLP}, 2018, pp. 2369--2380.

\bibitem{rajpurkar2018squad2}
P. Rajpurkar, R. Jia, and P. Liang, ``Know what you don't know: Unanswerable questions for SQuAD,'' in \textit{Proc. ACL}, 2018, pp. 784--789.

\bibitem{williams2018multinli}
A. Williams, N. Nangia, and S. Bowman, ``A broad-coverage challenge corpus for sentence understanding through inference,'' in \textit{Proc. NAACL-HLT}, 2018, pp. 1112--1122.

\bibitem{trivedi2022musique}
H. Trivedi, N. Balasubramanian, T. Khot, and A. Sabharwal, ``MuSiQue: Multihop questions via single hop question composition,'' \textit{Transactions of the ACL}, vol. 10, pp. 539--554, 2022.

\bibitem{schuster2021vitaminc}
T. Schuster, A. Fisch, and R. Barzilay, ``Get your vitamin C! Robust fact verification with contrastive evidence,'' in \textit{Proc. NAACL-HLT}, 2021, pp. 624--643.

\bibitem{zhang2015character}
X. Zhang, J. Zhao, and Y. LeCun, ``Character-level convolutional networks for text classification,'' in \textit{Proc. NeurIPS}, 2015, pp. 649--657.

\bibitem{saravia2018carer}
E. Saravia, H.-C. T. Liu, Y.-H. Huang, J. Wu, and Y.-S. Chen, ``CARER: Contextualized affect representations for emotion recognition,'' in \textit{Proc. EMNLP}, 2018, pp. 3687--3697.

\bibitem{robertson2009bm25}
S. Robertson and H. Zaragoza, ``The probabilistic relevance framework: BM25 and beyond,'' \textit{Foundations and Trends in Information Retrieval}, vol. 3, no. 4, pp. 333--389, 2009.

\bibitem{liu1989lbfgs}
D. C. Liu and J. Nocedal, ``On the limited memory BFGS method for large scale optimization,'' \textit{Mathematical Programming}, vol. 45, no. 1--3, pp. 503--528, 1989.

\bibitem{naeini2015ece}
M. P. Naeini, G. Cooper, and M. Hauskrecht, ``Obtaining well calibrated probabilities using Bayesian binning into quantiles,'' in \textit{Proc. AAAI}, 2015, pp. 2901--2907.

\bibitem{loshchilov2019decoupled}
I. Loshchilov and F. Hutter, ``Decoupled weight decay regularization,'' in \textit{Proc. ICLR}, 2019.

\bibitem{efron1993bootstrap}
B. Efron and R. J. Tibshirani, \textit{An Introduction to the Bootstrap}. New York: Chapman \& Hall, 1993.

\bibitem{mcnemar1947note}
Q. McNemar, ``Note on the sampling error of the difference between correlated proportions or percentages,'' \textit{Psychometrika}, vol. 12, no. 2, pp. 153--157, 1947.

\bibitem{dietterich1998approximate}
T. G. Dietterich, ``Approximate statistical tests for comparing supervised classification learning algorithms,'' \textit{Neural Computation}, vol. 10, no. 7, pp. 1895--1923, 1998.

\bibitem{zhao2021calibrate}
Z. Zhao, E. Wallace, S. Feng, D. Klein, and S. Singh, ``Calibrate before use: Improving few-shot performance of language models,'' in \textit{Proc. ICML}, 2021, pp. 12697--12706.

\bibitem{qwen2025qwen3}
Qwen Team, ``Qwen3 technical report,'' \textit{arXiv:2505.09388}, 2025.

\end{thebibliography}
\end{fullwidth}


\makeatletter
\if@twocolumn

  \begin{biography}[author1]{Subham Divakar} is a Senior Software Engineer at Soci\'{e}t\'{e} G\'{e}n\'{e}rale and a graduate of the Department of Electronics and Telecommunications, C.V. Raman Global University, Bhubaneswar, India. His research interests include attention mechanisms, retrieval-augmented generation, decision models for agentic AI, vision transformers, and deep learning for plant disease recognition. He has published IEEE papers on recursive attention mechanisms and plant disease classification. He also runs a social media marketing agency and maintains several open-source AI projects. (Email: shubham.divakar@gmail.com; ORCID: 0000-0003-0635-1058)
  \end{biography}

  \begin{biography}[author2]{Rojalina Priyadarshini} is affiliated with the Department of Computer Science, C.V. Raman Global University, Bhubaneswar, India. Her research interests include machine learning, natural language processing, information retrieval, and image processing. She has co-authored research on attention mechanisms, plant disease classification, and topology-aware AI systems. (Email: priyadarshini.rojalina@gmail.com; ORCID: 0000-0002-5481-5251)
  \end{biography}

\else
\fi
\makeatother

\end{document}
